\scanpage{644}
\mspage{32}

\uline{Corollaire} 2.11. Le prédérivateur $\unsure{\mathbb{D}|\underline{D}}_{M,W}$ défini
par un localiseur \uline{très spécial} $W$ dans $M$
est un dérivateur au plein sens du terme,
i.e. il satisfait à tous les axiomes
Der 1 à Der 6 (y inclus Der 4 bis et
Der 5 bis et Der 6 bis) des dérivateurs.

\uline{NB} À ma propre surprise, je n'ai pas été capable
de \unsure{trouver|donner} un cas d'un localiseur $W$,
qui définisse un dérivateur sur Cat tout
entier (avec existence des $u_!^W,u_*^W$),
\uline{que} : pour $W$ un localiseur très spécial,
donc \unsure{s'inscrivant|s'insérant} dans \unsure{:} un triple de
Quillen très spécial -- ou quand on
peut s'y ramener de la façon \unsure{standart|standard}.

(Par exemple dans le cas \deleted{\illegible{1}}\inserted{où} $M=\mathrm{Cat}$, $W$
localiseur fondamental \inserted{satisfaisant Loc \unsure{5|6} a) et Loc 8, 8 bis} -- les
mono $\in W$ se refusent formellement à
être stables par cochangement de base !)

% slanted
\margin{axiome des carrés cocartésiens (version directe), axiome des limites inductives, accessibilité}

J'ignore totalement si le résultat mirabolant
``prouvé'' par Anderson est vrai ou faux. \unsure{Ça|Sa} a l'air
\inserted{presque} trop beau pour être vrai \ldots. D'ailleurs, je ne
connais pas de structure de Quillen spéciale qui
ne soit très spéciale -- désolé !
% archive sheet number: 643
